\scanpage{69}
On pose
\[
\left\{\begin{aligned}
\mathrm{cn}^2z&=1-\mathrm{sn}^2z=\frac{\unsure{\wp|p}z-e_2}{\wp z-e_1}=\deleted{\unsure{-\wp(z-\omega_2)|-\wp(z-\omega_1)}}-\frac1{k^2}\bigl[\wp(z-\omega_1)-e_3\bigr]
&\boxed{\mathrm{cn}\,0=+1}\\
\mathrm{dn}^2z&=1-k^2\mathrm{sn}^2z=\frac{\wp z-e_3}{\wp z-e_1}=-\bigl[\wp(z-\omega_1)-e_2\bigr]
&\boxed{\mathrm{dn}\,0=+1}
\end{aligned}\right. \tag*{8}
\]

On a
\[
\left\{\begin{aligned}
\mathrm{cn}\,z&=\frac{\sigma\omega_1}{\sigma\omega_2}e^{(\eta_{2}-\eta_1)z}\frac{\sigma(z-\omega_2)}{\sigma(z-\omega_1)}
=\frac{\sigma\omega_1}{\sigma\omega_2}e^{-(\eta_2-\eta_1)z}\frac{\sigma(z+\omega_2)}{\sigma(z+\omega_1)}\\
\mathrm{dn}\,z&=\frac{\sigma\omega_1}{\sigma\omega_3}e^{(\eta_3-\eta_1)z}\frac{\sigma(z-\omega_3)}{\sigma(z-\omega_1)}
=\frac{\sigma\omega_1}{\sigma\omega_3}e^{-(\eta_3-\eta_1)z}\frac{\sigma(z+\omega_3)}{\sigma(z+\omega_1)}
\end{aligned}\right. \tag*{9}
\]

\[
\left\{\begin{aligned}
&{\times}\ \mathrm{cn}\,{-z}=\mathrm{cn}\,z &&\mathrm{dn}\,{-z}=\mathrm{dn}\,z &&\mathrm{cn}\,z\text{ et }\mathrm{dn}\,z\text{ sont paires}\\
&{\times}\left\{\begin{aligned}
\mathrm{cn}(z+2\omega_1)&=-\mathrm{cn}\,z\\
\mathrm{cn}(z+2\omega_2)&=-\mathrm{cn}\,z
\end{aligned}\right. &&
\begin{aligned}\mathrm{dn}(z+2\omega_3)&=-\mathrm{dn}\,z\\
\mathrm{dn}(z+2\omega_1)&=-\mathrm{dn}\,z\end{aligned}\\
&{\times}\left\{\begin{aligned}
\mathrm{cn}(z+2\omega_3)&=\mathrm{cn}\,z\\
\mathrm{cn}(z+4\omega_1)&=\mathrm{cn}\,z\\
\mathrm{cn}(z+4\omega_2)&=\mathrm{cn}\,z
\end{aligned}\right. &&
\begin{aligned}\mathrm{dn}(z+2\omega_2)&=\mathrm{dn}\,z\\
\mathrm{dn}(z+4\omega_3)&=\mathrm{dn}\,z\\
\mathrm{dn}(z+4\omega_1)&=\mathrm{dn}\,z\end{aligned} &&
\begin{aligned}&\mathrm{cn}\text{ admet pour périodes}\\&2\omega_3,4\omega_1,4\omega_2\\
&\mathrm{dn}\text{ admet pour période}\\&2\omega_3,4\omega_{\unsure{1|3}},4\omega_2\end{aligned}
\end{aligned}\right. \tag*{10}
\]

$\mathrm{cn}\,z$ et $\mathrm{dn}\,z$ sont fonctions elliptiques d'ordre 2 ;

\uuline{Poles de $\mathrm{cn}\,z$} homologues de $\omega_1$, résidu $\frac{i}{k}$ \inserted{si $\omega_1,\omega_2$ \deleted{est direct}}, et \deleted{$\omega_{\unsure{1|2}}$} \deleted{\unsure{(orienté) vers $2\omega_2$|(si $\omega_2$) vers $2\omega_2$}} \inserted{$-\omega_1$} résidu $-\frac{i}{k}$

(l'ambiguité de signe est levée par 9, joint à 6)

\deleted{Poles} \uuline{Zéros de $\mathrm{cn}\,z$} homologues de $+\omega_2$ et $-\omega_2$ ; on a,
\[
\mathrm{cn}'\omega_2=-\frac{e^{\eta_2\omega_2}}{\sigma^2\omega_2}\qquad (9,\text{ puis }6)
\]
; a partir de relation de définition de $\mathrm{cn}\,z$, on montre $\mathrm{cn}^{\prime\,2}\omega_2=-\mathrm{sn}''\omega_2$
\[
\mathrm{cn}'\omega_2=\pm\sqrt{1-k^2}
\]
Enfin
\[
\mathrm{cn}\,\omega_3=\frac{i}{k}\frac{e^{\eta_2\omega_2}}{\sigma^2\omega_2}=-\frac{i}{k}\mathrm{cn}'\omega_2=\pm\frac{i}{k}\sqrt{1-k^2}\qquad \mathrm{cn}'\omega_3=0
\]

% archive sheet number: 68
